% !TEX program = xelatex
% AI-effects 프로젝트 --- 점유율 지표 A·B·C와 노출도 3개의 순위 격차 상위·하위 30 엑셀 파일 설명.
% 표 본문: results/table/36_summary.tex, 36_robust_pos.tex, 36_robust_neg.tex, 36_notes.tex
%          (code/04_19_analysis_kor_share_exposure_rankgap.do가 생성)
% 설명 대상: results/table/36_kor_share_exposure_rankgap.xlsx
% 관련: paper 28(점유율과 노출도 상관), 29(점유율과 노출도의 격차), 35(순위 비교 목록)
% 컴파일: xelatex -> biber -> xelatex -> xelatex (kotex, biblatex/biber 필요)
\documentclass[11pt]{article}

\usepackage{fontspec}
\usepackage{kotex}
\usepackage[a4paper,margin=2.0cm]{geometry}
\usepackage{longtable}
\usepackage{booktabs}
\usepackage{array}
\usepackage{amsmath}
\usepackage{xcolor}
\usepackage{hyperref}
\usepackage{enumitem}
\usepackage[backend=biber,style=authoryear,maxcitenames=2,uniquelist=false,sorting=nyt]{biblatex}
\addbibresource{../../reference/references.bib}

\hypersetup{colorlinks=true,linkcolor=blue!50!black,urlcolor=blue!50!black}
\newcolumntype{L}[1]{>{\raggedright\arraybackslash}p{#1}}
\newcolumntype{R}[1]{>{\raggedleft\arraybackslash}p{#1}}
\setlength{\LTleft}{0pt}
\setlength{\LTright}{0pt}
\renewcommand{\arraystretch}{1.15}
\setlength{\tabcolsep}{3.5pt}

\title{AI 점유율 지표와 노출도의 순위 격차 엑셀 파일 설명\\[2mm]
\large --- 9개 쌍(지표 A$\cdot$B$\cdot$C $\times$ Felten$\cdot$Eloundou$\cdot$Observed)의 격차 상위$\cdot$하위 30 ---}
\author{AI-effects 프로젝트}
\date{\today}

\begin{document}
\maketitle

\begin{abstract}
\noindent 이 노트는 \texttt{results/table/36\_kor\_share\_exposure\_rankgap.xlsx}를 설명한다. 파일은 \texttt{code/04\_19\_analysis\_kor\_share\_exposure\_rankgap.do}가 만든다. 점유율 지표 A$\cdot$B$\cdot$C(실제 사용)와 노출도 3개(Felten AIOE\parencite{felten-2023-how-will-language-modelers-like-chatgpt}, Eloundou $\beta$\parencite{eloundou-2023-gpts-are-gpts-an-early}, Observed Exposure\parencite{massenkoff-2026-labor-market-impacts-of-ai})를 짝지은 9개 쌍마다, KSCO 8차 소분류 129개의 순위 격차가 가장 큰 30개(사용 $>$ 노출도)와 가장 작은 30개(사용 $<$ 노출도)를 시트로 담았다.
\begin{itemize}[leftmargin=*]
  \item \textbf{순위 격차} $G = r(\text{노출도}) - r(\text{사용})$. 양수이면 사용 순위가 노출도 순위보다 높다. 순위는 paper 35와 같다.
  \item \textbf{시트 19개}: 쌍별 격차 상위 30 9개(\texttt{*\_pos}), 하위 30 9개(\texttt{*\_neg}), 129개 전부의 순위와 격차 9개(\texttt{all}).
  \item \textbf{격차는 크다.} 쌍마다 $|G|$ 평균은 22$\sim$27계단이고, 30계단 이상 벌어진 소분류가 33$\sim$48개다. Felten과의 격차가 가장 크다.
  \item \textbf{쌍을 바꿔도 같은 직업이 나온다.} 9개 쌍 모두의 상위 30에 드는 소분류가 11개, 하위 30에 드는 소분류가 13개다. 사용이 노출도보다 앞서는 쪽은 연구$\cdot$공학$\cdot$영상 전문가와 기능$\cdot$설비 직업이고, 노출도가 앞서는 쪽은 보건(영양사, 약사, 간호사), 판매$\cdot$상담, 사무 직업이다.
\end{itemize}
\end{abstract}

\section{파일 구성}

\begin{itemize}[leftmargin=*]
  \item \textbf{쌍별 시트}(18개): 이름은 \texttt{\{사용 지표\}\_\{노출도\}\_\{pos$|$neg\}}다. 사용 지표는 \texttt{A}, \texttt{B}, \texttt{C}, 노출도는 \texttt{FEL}(Felten AIOE), \texttt{ELO}(Eloundou $\beta$), \texttt{OE}(Observed Exposure)다. \texttt{pos}는 격차 상위 30(사용 $>$ 노출도), \texttt{neg}는 격차 하위 30(사용 $<$ 노출도)이다. 요청 목록 (8)$\sim$(16)이 \texttt{pos}, (17)$\sim$(25)가 \texttt{neg}다.
  \item \textbf{쌍별 시트의 열}: 격차 순위, KSCO 코드, 직업명, 사용 지표 값, 사용 지표 순위, 노출도 값, 노출도 순위, 순위 격차. 첫 열(격차 순위)은 요청 목록에 없던 열로, \texttt{pos}는 격차가 큰 순, \texttt{neg}는 격차가 작은(가장 음수인) 순이다.
  \item \textbf{참고 시트} \texttt{all}: 129개 소분류의 지표 A$\cdot$B$\cdot$C와 노출도 3개의 값$\cdot$순위, 9개 쌍의 순위 격차. 30위 밖의 격차나 다른 기준으로 다시 고를 때 쓴다.
\end{itemize}

\paragraph{정의.}
\begin{itemize}[leftmargin=*]
  \item 표본: 지표 A$\cdot$B$\cdot$C와 노출도 3개가 모두 있는 소분류 129개(paper 27$\cdot$28$\cdot$35와 같다).
  \item 순위 $r(\cdot)$: 129개 안에서 값이 큰 순서(1 = 가장 큼), 같은 값은 소분류 코드 순. 엑셀 35(paper 35)의 순위와 같다.
  \item 순위 격차 $G_{j}(m, x) = r_j(x) - r_j(m)$. $m$은 사용 지표(A, B, C), $x$는 노출도다. $G = +50$이면 그 직업의 사용 순위가 노출도 순위보다 50계단 높다. 129개 안에서 $\sum_j G_j = 0$이다.
  \item 상위 30은 $G$ 내림차순, 하위 30은 $G$ 오름차순으로 고른다. $G$가 같으면 코드 순이다.
\end{itemize}
지표와 노출도의 내용은 paper 35의 2절과 같다. 지표 A는 업무용 대화 점유율, B는 절감 시간 점유율, C는 B를 1인당 근로시간으로 나눈 취업자 환산 점유율이다.

\section{결과 개관}

\subsection{쌍별 요약}

\begin{small}
\begin{longtable}{L{1.3cm}L{2.6cm}R{1.0cm}R{1.3cm}R{1.4cm}R{1.4cm}R{1.2cm}R{1.2cm}R{1.5cm}R{1.5cm}}
\caption{쌍별 순위 격차 요약(소분류 129개)}\label{tab:summary}\\
\toprule
사용 & 노출도 & $\rho$ & $|G|$ 평균 & $|G| \ge 30$ & $G > 0$ & 최대 & 최소 & 상위 30 경계 & 하위 30 경계 \\
\midrule
\endfirsthead
% 04_19_analysis_kor_share_exposure_rankgap.do가 생성. 직접 고치지 말 것.
지표 A & Felten AIOE & 0.60 & 26.5 & 48 & 64 & +80 & -96 & +26 & -23 \\
지표 A & Eloundou β & 0.69 & 22.4 & 33 & 68 & +76 & -85 & +19 & -20 \\
지표 A & Observed Exposure & 0.70 & 22.2 & 38 & 64 & +88 & -78 & +20 & -19 \\
\midrule
지표 B & Felten AIOE & 0.63 & 25.4 & 45 & 65 & +77 & -95 & +25 & -22 \\
지표 B & Eloundou β & 0.70 & 22.0 & 35 & 66 & +80 & -83 & +21 & -18 \\
지표 B & Observed Exposure & 0.69 & 22.4 & 39 & 70 & +92 & -72 & +20 & -20 \\
\midrule
지표 C & Felten AIOE & 0.63 & 25.3 & 43 & 65 & +84 & -94 & +25 & -21 \\
지표 C & Eloundou β & 0.70 & 22.2 & 36 & 68 & +78 & -82 & +19 & -18 \\
지표 C & Observed Exposure & 0.69 & 22.5 & 39 & 68 & +90 & -73 & +21 & -20 \\

\bottomrule
\end{longtable}
\end{small}
\noindent $\rho$는 Spearman 순위상관, ``경계''는 상위$\cdot$하위 30번째 소분류의 격차다.

순위상관이 낮을수록 격차가 크다. Felten 쌍은 $|G|$ 평균이 25$\sim$27계단, $|G| \ge 30$인 소분류가 43$\sim$48개로 가장 크고, Eloundou$\cdot$Observed 쌍은 22계단 안팎이다. 상위 30에 들려면 격차가 +19$\sim$+26계단, 하위 30에 들려면 $-18\sim-23$계단 이상이어야 한다. 격차가 양수인 소분류(64$\sim$70개)가 음수보다 조금 많지만, 음수 쪽 최솟값($-72\sim-96$)이 더 깊은 쌍이 많다.

\subsection{여러 쌍에서 반복되는 소분류}

9개 쌍의 목록은 많이 겹친다. 상위 30 목록 9개를 합치면 서로 다른 소분류가 56개, 하위 30 목록은 52개뿐이다(표~\ref{tab:notes}). 아래 표는 9개 중 7개 이상의 목록에 든 소분류다. 순위는 129개 안 순위이고, 평균 격차는 그 소분류가 든 목록들의 평균이다.

\begin{small}
\begin{longtable}{L{5.6cm}R{1.2cm}R{1.4cm}R{1.2cm}R{1.2cm}R{1.3cm}R{1.4cm}R{1.4cm}}
\caption{격차 상위 30(사용 $>$ 노출도) 목록 9개 중 7개 이상에 든 소분류}\label{tab:robpos}\\
\toprule
소분류 & 목록 수 & 평균 $G$ & $r$(A) & $r$(C) & $r$(Felten) & $r$(Eloundou) & $r$(Observed) \\
\midrule
\endfirsthead
\toprule
소분류 & 목록 수 & 평균 $G$ & $r$(A) & $r$(C) & $r$(Felten) & $r$(Eloundou) & $r$(Observed) \\
\midrule
\endhead
% 04_19_analysis_kor_share_exposure_rankgap.do가 생성. 직접 고치지 말 것.
852 냉난방 관련 설비 조작원 & 9 & +88.8 & 37 & 35 & 101 & 121 & 149.5 \\
799 기타 기능 관련 종사자 & 9 & +65.3 & 32 & 28 & 119 & 108 & 64 \\
999 기타 서비스 관련 단순 종사자 & 9 & +61.0 & 52 & 76 & 133 & 144 & 116 \\
910 건설 및 광업 단순 종사자 & 9 & +60.0 & 91 & 92 & 163 & 161 & 133 \\
851 금속 공작 기계 조작원 & 9 & +56.0 & 97 & 86 & 149 & 151 & 134 \\
782 건설 관련 기능 종사자 & 9 & +55.0 & 78 & 83 & 152 & 140 & 114 \\
730 목재‧가구‧악기 및 간판 관련 기능 종사자 & 9 & +54.7 & 56 & 54 & 116 & 117 & 95 \\
721 섬유 및 가죽 관련 기능 종사자 & 9 & +54.3 & 69 & 77 & 132 & 123 & 130 \\
293 연극‧영화 및 영상 전문가 & 9 & +53.7 & 11 & 9 & 67 & 66 & 58 \\
213 생명 및 자연과학 관련 시험원 & 9 & +51.3 & 8 & 7 & 84 & 53 & 40 \\
149 기타 건설‧전기 및 생산 관련 관리자 & 8 & +57.9 & 46 & 34 & 81 & 69 & 126 \\
941 청소 관련 종사자 & 8 & +46.6 & 99 & 80 & 142 & 131 & 122 \\
235 기계‧로봇공학 기술자 및 시험원 & 8 & +44.3 & 15 & 17 & 56 & 41 & 78 \\
239 기타 공학 전문가 및 관련 종사자 & 8 & +42.1 & 23 & 21 & 65 & 46 & 73 \\
246 치료‧재활사 및 의료기사 & 7 & +44.9 & 38 & 47 & 71 & 84 & 99 \\
264 문리‧기술 및 예능 강사 & 7 & +38.6 & 14 & 11 & 39 & 61 & 45 \\

\bottomrule
\end{longtable}

\begin{longtable}{L{5.6cm}R{1.2cm}R{1.4cm}R{1.2cm}R{1.2cm}R{1.3cm}R{1.4cm}R{1.4cm}}
\caption{격차 하위 30(사용 $<$ 노출도) 목록 9개 중 7개 이상에 든 소분류}\label{tab:robneg}\\
\toprule
소분류 & 목록 수 & 평균 $G$ & $r$(A) & $r$(C) & $r$(Felten) & $r$(Eloundou) & $r$(Observed) \\
\midrule
\endfirsthead
\toprule
소분류 & 목록 수 & 평균 $G$ & $r$(A) & $r$(C) & $r$(Felten) & $r$(Eloundou) & $r$(Observed) \\
\midrule
\endhead
% 04_19_analysis_kor_share_exposure_rankgap.do가 생성. 직접 고치지 말 것.
245 영양사 & 9 & -80.0 & 103 & 98 & 7 & 19 & 37 \\
322 운송 및 무역 사무원 & 9 & -69.7 & 98 & 85 & 26 & 13 & 20 \\
531 통신 관련 판매 종사자 & 9 & -61.0 & 60 & 72 & 1 & 8 & 12 \\
243 약사 및 한약사 & 9 & -56.7 & 104 & 111 & 50 & 55 & 50 \\
311 기획 및 마케팅 사무원 & 9 & -54.7 & 63 & 84 & 35 & 2 & 27 \\
159 기타 판매 및 고객 서비스 관리자 & 9 & -45.0 & 86 & 71 & 30 & 27 & 33 \\
244 간호사 & 9 & -41.3 & 106 & 108 & 69 & 67 & 61 \\
362 고객 상담 및 모니터 요원 & 9 & -40.7 & 44 & 53 & 18 & 6 & 4 \\
134 문화‧예술 관련 관리자 & 9 & -38.0 & 89 & 79 & 41 & 33 & 58 \\
238 항공기‧선박 기관사 및 관제사 & 9 & -36.0 & 118 & 113 & 82 & 71 & 83 \\
312 인사 및 교육·훈련 사무원 & 9 & -35.0 & 43 & 46 & 8 & 15 & 7 \\
451 조리사 & 9 & -29.3 & 129 & 129 & 95 & 106 & 98 \\
429 기타 돌봄 및 보건 서비스 종사자 & 9 & -26.7 & 119 & 116 & 87 & 95 & 91 \\
331 회계 및 경리 사무원 & 8 & -38.6 & 45 & 57 & 33 & 4 & 17 \\
532 방문 및 노점 판매 종사자 & 7 & -67.4 & 81 & 104 & 10 & 85 & 31 \\
441 운송 서비스 종사자 & 7 & -41.3 & 101 & 112 & 84 & 89 & 46 \\
423 노인 및 장애인 돌봄 종사자 & 7 & -34.7 & 115 & 110 & 81 & 94 & 69 \\
112 기업 대표 및 기업 고위 임원 & 7 & -31.0 & 72 & 59 & 31 & 29 & 53 \\

\bottomrule
\end{longtable}
\end{small}

\paragraph{사용이 노출도보다 앞서는 직업.} 주로 두 묶음이다.
\begin{itemize}[leftmargin=*]
  \item \textbf{연구$\cdot$공학$\cdot$영상 전문가}(213, 211, 293, 235, 239, 264). 사용 순위는 7$\sim$23위로 높은데 노출도 순위는 31$\sim$83위다. 노출도가 중간인 직업에서 실제 대화가 많이 나온 경우다.
  \item \textbf{기능$\cdot$설비$\cdot$단순 직업}(852, 799, 730, 721, 782, 851, 910, 999, 941). 노출도는 대부분 100위 밖으로 낮은데(Observed는 62$\sim$125위) 사용은 28$\sim$99위의 중하위권이다. 노출도 순위가 바닥이라 사용이 조금만 있어도 격차가 크게 나온다. 852 냉난방 관련 설비 조작원은 Observed Exposure가 0이라 그 순위(125위)가 동점 처리의 결과다.
\end{itemize}

\paragraph{노출도가 사용보다 앞서는 직업.} 세 묶음이다.
\begin{itemize}[leftmargin=*]
  \item \textbf{보건}(245 영양사, 243 약사 및 한약사, 244 간호사, 429 기타 돌봄 및 보건 서비스). 노출도 순위는 7$\sim$95위인데 사용 순위는 98$\sim$119위다. 245 영양사는 Felten 7위, 지표 A 103위로 평균 격차가 $-80$이다.
  \item \textbf{판매$\cdot$상담$\cdot$관리}(531 통신 관련 판매, 159 기타 판매 및 고객 서비스 관리자, 362 고객 상담 및 모니터 요원, 532 방문 및 노점 판매, 134 문화$\cdot$예술 관련 관리자). Felten에서 특히 높다(531은 Felten 1위).
  \item \textbf{사무}(322 운송 및 무역 사무원, 311 기획 및 마케팅 사무원, 312 인사 및 교육$\cdot$훈련 사무원, 331 회계 및 경리 사무원). 세 노출도 모두 2$\sim$35위로 높지만 사용 순위는 43$\sim$98위다. paper 29에서 ``노출도 $>$ 사용''으로 나온 311과 331이 여기에도 든다.
\end{itemize}

\section{읽을 때 유의사항}

\begin{longtable}{L{12.5cm}R{2.5cm}}
\caption{점검 수치}\label{tab:notes}\\
\toprule
항목 & 값 \\
\midrule
\endfirsthead
% 04_19_analysis_kor_share_exposure_rankgap.do가 생성. 직접 고치지 말 것.
비교 소분류 수 / 그중 사용 0(지표 A·B·C = 0) & 163 / 34 \\
목록 칸 수: 상위 30 목록 / 하위 30 목록(9개 쌍 합, 경계 동점 포함) & 277 / 272 \\
사용 0인 소분류가 든 칸: 상위 30 목록 / 하위 30 목록 & 0 / 78 \\
9개 쌍 상위 30 목록에 한 번 이상 든 소분류 수 / 9개 모두에 든 수 & 60 / 10 \\
9개 쌍 하위 30 목록에 한 번 이상 든 소분류 수 / 9개 모두에 든 수 & 55 / 10 \\
Observed Exposure = 0인 소분류 수 & 28 \\
Observed Exposure 쌍 상위 30(3개 목록) 중 Observed Exposure = 0인 칸 & 23 / 93 \\
한 달만 공개된 소분류 수 & 7 \\
한 달만 공개된 소분류가 든 칸: 상위 30 목록 / 하위 30 목록 & 6 / 3 \\

\bottomrule
\end{longtable}

\begin{enumerate}[leftmargin=*]
  \item \textbf{순위 격차는 분포의 끝에서 작아진다.} 사용 1위 직업은 양의 격차를 가질 수 없고, 129위 직업은 음의 격차를 가질 수 없다. 그래서 사용 최상위 직업(222, 291 등)은 노출도도 높아 목록에 잘 들지 않고, 목록은 중간 순위 직업 위주가 된다. 강도 자체(몇 배 차이인지)는 paper 29의 로그 비율 격차로 본다.
  \item \textbf{Observed Exposure가 0인 소분류 13개는 117$\sim$129위 동점이다.} 코드 순으로 순위를 매겼기 때문에 이 소분류의 Observed 쌍 격차는 최대 12계단까지 코드 순서에 따라 달라진다. Observed 쌍 상위 30 목록 90칸 중 12칸이 이 소분류다(852, 710, 881, 882, 422, 453).
  \item \textbf{한 달만 공개된 소분류는 하위 30 쪽에만 든다.} 7개 중 사용 순위가 바닥인 소분류(451 조리사 등)가 하위 30 목록 14칸에 든다. 비공개 달을 0으로 두어 사용 값이 절반이 된 영향이 섞여 있다(paper 27).
  \item \textbf{점유율은 직업 규모가 들어간 값이다.} 큰 직업은 점유율 순위가 오르기 쉬워 격차가 양수 쪽으로 간다. 규모를 뺀 비교는 상대 활용도(paper 25, 32)나 노출도 기반 점유율(paper 28, 33)로 본다.
  \item \textbf{두 쪽 모두 SOC 값을 연계표로 옮긴 것이다}(paper 21, 23). 한 SOC에 쏠린 소분류는 그 SOC의 값에 따라 격차가 크게 움직인다. 격차가 큰 직업은 SOC 구성(paper 26)을 함께 확인하는 것이 좋다.
\end{enumerate}

\section{산출 파일}

\begin{itemize}[leftmargin=*]
  \item 코드: \texttt{code/04\_19\_analysis\_kor\_share\_exposure\_rankgap.do}(Stata 17, \texttt{00\_master.do}에 포함). 입력은 \texttt{data/proc/merged/kor\_ai\_usage\_ksco3.dta}(\texttt{03\_05}).
  \item 엑셀: \texttt{results/table/36\_kor\_share\_exposure\_rankgap.xlsx}(시트 19개).
  \item 이 노트의 표: \texttt{results/table/36\_summary.tex}, \texttt{36\_robust\_pos.tex}, \texttt{36\_robust\_neg.tex}, \texttt{36\_notes.tex}.
\end{itemize}

\printbibliography[title={참고문헌}]

\end{document}
